\documentclass[10pt,a4paper]{amsart}
\usepackage[margin=19mm]{geometry}
\usepackage{amsmath,amssymb,mathtools}
\usepackage[scr=rsfs,cal=euler]{mathalfa}
\usepackage{xfrac}
\usepackage[unicode,hidelinks,bookmarks=false]{hyperref}
\newcommand{\ie}{i.e.\ }
\newcommand{\cf}{cf.\ }
\newcommand{\resp}{respectively\ }
\newcommand{\Subsection}{\subsection}
\providecommand{\nobreakdash}{\nobreak\hbox{-}}
\setlength{\emergencystretch}{2em}
\multlinegap0pt
\newtheorem{cor}{Corollary}
\newtheorem{prop}{Proposition}
\renewcommand{\a}[0]{{\mathbb A}} 
\newcommand{\hooklongrightarrow}{\lhook\joinrel\longrightarrow}
\renewcommand{\o}[0]{{\mathcal O}} 
\newcommand{\p}[0]{{\mathbb P}}
\renewcommand{\r}[0]{{\mathbb R}} 
\title{Resonance of rank-two vector bundles over elliptic curves}
\author{C\u{a}lin Spiridon}
\date{Source: arXiv:2604.23291v1 (CC BY 4.0)}
\begin{document}
\maketitle

\noindent\emph{Source and licence.} Resonance of rank-two vector bundles over elliptic curves. Version arXiv:2604.23291v1 (CC BY 4.0), distributed by arXiv under the Creative Commons Attribution 4.0 International licence, http://creativecommons.org/licenses/by/4.0/, as recorded in the arXiv licence metadata for this exact version and shown on the abstract page https://arxiv.org/abs/2604.23291v1. Full licence notice: \texttt{licenses/arxiv-2604.23291-CC-BY-4.0.txt}.\par\smallskip

\noindent\textit{Editorial scope.} Counted unit: the article's prop prop:2P\_bQ (source0.tex lines 318--320). The article's complete original proof of it is delivered with it (source0.tex lines 322--334, one \texttt{proof} environment of 13 lines). No mathematics is written, repaired or re-derived: the statement and the proof are the article's own lines, in the article's own notation, and no retained line is substituted or re-flowed.  Retained uncounted as the source-local context the proof needs: lines 256--265.  Nothing was omitted from any retained span, so the package carries no omission record.  No retained line contains a citation, so the wrapper carries no bibliography and no bibliography entry.  No retained line contains a figure environment, an image-inclusion command or drawing code, so no image is delivered and the package declares no asset dependency.  Editorial formatting only, in addition: an AMS \texttt{amsart} preamble that restates the article's own declarations and macros used by the retained lines, namely the environments \texttt{cor}, \texttt{prop} on separate counters, together with the standard packages those lines need.  The retained environments print in this document as Corollary 1, Proposition 1 in order of appearance, whereas the article prints the same items with the numbering of its own full document.  The complete native source of the article as distributed in the arXiv e-print for this exact version is bundled unchanged as \texttt{sources/199151/source0.tex}.
\begin{cor} \label{cor:saturated_embeddings_a<b}
If $E = \o_C(aP) \oplus \o_C(bQ)$ with $a<b$ and $L = \o_C(dR)$ is a degree $d$ line bundle on $C$, then
\begin{itemize}
    \item [(a)] If $d > b$, then $L$ cannot be embedded into $E$.
    \item [(b)] If $d = b$, then $L$ can be embedded into $E$ if and only if $bR \sim bQ$, that is $L = \o_C(bQ)$. In this case, $L$ embeds uniquely (up to scalar multiplication) and saturated in $E$.
    \item [(c)] If $a < d < b$, then $L$ cannot be embedded into $E$ as a saturated sub-line bundle.
    \item [(d)] If $d = a$, then $L$ can be embedded into $E$ as a saturated sub-line bundle if and only if $aR \sim aP$, that is $L = \o_C(aP)$. In this case, any saturated embedding of $L$ into $E$ is determined by a section of $H^0(\o_C(bQ - aP))$.
    \item [(e)] If $d < a$, then $L$ can be embedded into $E$ as a saturated sub-line bundle. Moreover, any saturated embedding of $L$ in $E$ is determined by a pair of two non-zero sections $(s_a, s_b) \in H^0(\o_C(aP - dR)) \oplus H^0(\o_C(bQ - dR))$ without common zeros. 
\end{itemize}
\end{cor}

\begin{prop} \label{prop:2P_bQ}
Let $E = \o_C(2P) \oplus \o_C(bQ)$, with $b \ge 3$. Then $\overline{\r_2(E)}$ and $\r_b(E)$ are the irreducible components of the resonance of $E$.
\end{prop}

\begin{proof}
Due to Corollary \ref{cor:saturated_embeddings_a<b}, the only line bundles that can be embedded as saturated subpencils into $E$ are $L = \o_C(2P)$ and $M = \o_C(bQ)$. Therefore, the degree stratification of the resonance consists of two strata. The stratum $\r_2(E)$ is simply the fiber of $\rho_2$ over $L = \o_C(2P)$, which is the image of
\[
\theta_L : \p H^0(\o_C(2P)) \times U_L \longrightarrow \p H^0(E),
\]
where the locus $U_L \subseteq \p H^0(\o_C \oplus \o_C(bQ - 2P))$ of saturated embeddings of $L$ into $E$ may be identified with $H^0(\o_C(bQ - 2P))$. It can be readily seen that the map $\theta_L$ is injective and hence, the stratum $\r_2(E)$ is isomorphic to $\p^1 \times \a^{b-2} \hooklongrightarrow \p^{b+1}$.

The closed stratum $\r_b(E)$ is the fiber of $\rho_b$ over $M = \o_C(bQ)$, which is simply the image of the map
\[
\theta_M : \p H^0(\o_C(bQ)) \hooklongrightarrow \p H^0(E). 
\]
Therefore, $\r_b(E)$ identifies with $\p^{b-1} \hooklongrightarrow \p^{b+1}$. Since $\overline{\r_2(E)}$ and $\r_b(E)$ are irreducible of the same dimension, we infer that these are the irreducible components of $\r(E)$.
\end{proof}

\end{document}
